\documentclass[11pt,a4paper]{article}

\usepackage{amsmath,amssymb,amsthm}
\usepackage{mathtools}
\usepackage{enumitem}
\usepackage{booktabs}
\usepackage{hyperref}
\hypersetup{
    colorlinks=true,
    linkcolor=blue,
    citecolor=blue,
    urlcolor=blue,
    pdfauthor={Anonymous},
    pdftitle={Cred Part 3 (Sketch): Graded Proofs, Zero Evidence, and Many-Valued Metatheory}
}
\usepackage{cleveref}
\usepackage[margin=1in]{geometry}

\newtheorem{theorem}{Theorem}[section]
\newtheorem{definition}[theorem]{Definition}
\newtheorem{remark}[theorem]{Remark}

\newcommand{\cred}{\mathsf{cred}}
\newcommand{\Cred}{\mathsf{Cred}}
\newcommand{\negc}{\mathord{\sim}}
\newcommand{\conjc}{\otimes}
\newcommand{\disjc}{\sqcup}
\newcommand{\condc}[2]{({#1} \mid {#2})}
\newcommand{\Real}{\mathbb{R}}

\title{\textbf{Cred Part~3 (Sketch): Graded Proofs, Zero Evidence, and Many-Valued Metatheory}}
\author{Anonymous (working draft)}
\date{February 2026}

\begin{document}
\maketitle

\begin{abstract}
Parts~1--2 present \emph{Cred} as a real-valued credence algebra on $[0,1]$ with
primitive conditioning governed by the chain rule
$\cred(A \mid B) \cdot \cred(B) = \cred(A \land B)$, deliberately leaving the
evidence-zero case underdetermined.  Part~3 is the ``how to use this as a
foundation'' layer: (i) graded proof patterns and dynamics (asymptotic proof,
robust proof); (ii) a careful interface to the existing literature on
conditional probability as primitive, conditional events, coherence, and
imprecise probability (especially conditioning on probability-zero events);
and (iii) a map from Cred's logic-style instantiations to established many-valued
logics and many-valued mathematics (product logic, BL/MTL families, and
many-valued set theory), with self-reference handled by fixed points rather than
contradiction.

This document is a \emph{paper sketch}: a structured synthesis and literature
map that is intended to stabilize the Part~3 roadmap and prevent re-deriving
known results under new names.
\end{abstract}

\section{What Part~3 Is For}

Part~1 gives the core algebra (values, complement, product, De~Morgan dual, and
primitive conditioning).  Part~2 supplies the missing layers needed to do
reasoning (valuations, consequence relations, update rules, and graded
predicates/quantifiers).  Part~3 addresses what remains once the primitives are
fixed:
\begin{enumerate}[label=(\roman*),nosep]
\item how ``proof'' and ``inference'' behave when credences are graded and
changing over time;
\item what to do at the evidence-zero boundary (conditioning on null events);
\item how Cred relates to the existing many-valued and probabilistic foundations
literature, so that we do not claim novelty where there is none.
\end{enumerate}

\section{Two Views, One Bridge}

Cred supports two complementary uses.

\subsection{Probability-style (joints are external)}

In the probability-style view, a valuation supplies marginals and \emph{joints}
as external data (from empirical models, domain constraints, or higher-level
structure).  Conditioning is constrained by the chain rule and is unique when
the evidence credence is positive.

\subsection{Logic-style (fix a truth-functional joint)}

In the logic-style view, one chooses a \emph{fixed} truth-functional joint
$j(a,b)$ and thereby induces a binary ``arrow'' (a truth-functional conditional)
for positive evidence:
\[
f(a,b) := \frac{j(a,b)}{a}\qquad (a>0).
\]
When one takes the Fr\'echet upper bound (maximal dependence) $j(a,b)=\min(a,b)$,
then for $a>0$ one obtains
\[
f(a,b)=\min(b/a,1),
\]
which is exactly the \emph{product-residuated (Goguen) implication} familiar
from product fuzzy logic~\cite{hajek1998,goguen1969}.
Cred's divergence from product logic is the evidence-zero boundary: product
logic sets $f(0,b)=1$ (ex falso/vacuous truth), while Cred treats evidence-zero
conditioning as underdetermined by design.

\subsection{Why the evidence-zero row is special}

The chain rule is well-formed at evidence zero and forces only the joint
to be zero; it does \emph{not} choose a specific conditional value.  This is
closely related to the classical ambiguity of ratios $j/a$ as $(j,a)\to (0,0)$:
the limiting value depends on the relative rate at which $j$ decays compared to
$a$, and no single value is canonical.

This boundary is also the main point of contact with established work on
conditioning on null events (measure-theoretic disintegration and regular
conditional probabilities) and with coherence-based conditional probability and
conditional events.

\section{Many-Valued Logics and ``Many-Valued Mathematics''}

Many-valued logics are not just alternative truth tables; they come with a large
algebraic and proof-theoretic ecosystem (residuated lattices, BL/MTL families,
Pavelka-style graded provability, and applications to fuzzy sets and topology).
The right synthesis question for Cred is not whether product logic exists (it
does), but what Cred adds, and which pieces it declines to commit to.

\subsection{Mathematical fuzzy logic (product logic in context)}

H\'ajek's \emph{metamathematics of fuzzy logic} organizes the ``big three''
continuous t-norm logics: \L ukasiewicz, G\"odel, and product~\cite{hajek1998}.
The corresponding algebraic semantics are BL-algebras and their special cases;
see also standard references on t-norms and residuated structures
in fuzzy logic~\cite{klement2000,bergmann2008,cintula2011,galatos2007}.

An important proof-theoretic strand is \emph{Pavelka-style} fuzzy logic, where
provability itself is graded and completeness is formulated in terms of degrees.
Pavelka's 1979 series develops many-valued rules of inference, semantics via
enriched residuated lattices, and semantic completeness for many-valued
calculi~\cite{pavelka1979i,pavelka1979ii,pavelka1979iii}.  This literature is a
useful baseline for Part~3: whenever we talk about graded consequence or graded
proof, we should be explicit about whether we mean (a) a semantic threshold
consequence relation (as in Part~2), or (b) a proof system whose theorems come
with provability degrees (Pavelka-style).

Product logic is a fully developed many-valued logic:
\begin{itemize}[nosep]
\item it is $[0,1]$-valued;
\item its implication is the product residuum (Goguen implication);
\item it validates a graded form of modus ponens via the adjointness
relationship $a\conjc f(a,b)\le b$.
\end{itemize}

Cred's logic-style instantiation with $j(a,b)=\min(a,b)$ agrees with product
logic \emph{for positive evidence} and then intentionally refuses to collapse
the evidence-zero boundary into a single forced value.

\subsection{Many-valued set theory and graded foundations}

There is also a tradition of ``many-valued mathematics'' in which set-theoretic
or logical foundations are valued in $[0,1]$ or other lattices.  For example,
Takeuti and Titani develop a many-valued set theory based on fuzzy logic and
study standard set-theoretic issues in that setting~\cite{takeuti1984}.
This literature provides concrete guidance on which foundational questions are
substantive (e.g., extensionality, comprehension, and interpretation of
membership), and which are merely notational changes.

\section{Conditional Probability as Primitive and Conditional Events}

Cred's conditioning is directly aligned with two established traditions:

\subsection{Primitive conditional probability}

R\'enyi and Popper develop axiom systems where conditional probability is
primitive and the chain rule is central~\cite{renyi1955,popper1959}.
van~Fraassen develops and applies Popper-style approaches in the analysis of
probabilities of conditionals and in the connection between probability and
relevant logic~\cite{vanfraassen1976,vanfraassen1983}.

\subsection{Conditional events and coherence}

De~Finetti treats conditional events as three-valued (true/false/void) in a
betting/coherence setting, providing an explicit semantics for the antecedent-false
case that avoids vacuous truth~\cite{definetti1936,definetti1937}.
Calabrese formalizes conditional events algebraically as Boolean fractions and
analyzes their operations~\cite{calabrese1987}.
The modern logic of indicative conditionals has also revisited de~Finetti-style
tables and their proof theory and algebraic semantics; see, e.g., the
De~Finettian logics DF/TT and related systems studied by \'Egr\'e, Rossi, and
Sprenger~\cite{egre2021}.
More recently, the coherence-based literature develops conditional random
quantities and operations on conditionals, including conjunction and iterated
conditioning~\cite{gilio2014,gilio2013}.  Coletti and Scozzafava provide a broad
synthesis of probabilistic logic in a coherence setting and devote substantial
attention to probability-zero events and conditionalization at the boundary
case~\cite{coletti2002}.

From the Cred perspective, these works are especially relevant for Part~3
because they show how to reason coherently when antecedents have probability
zero \emph{without} forcing a single truth-functional value, and because they
explicitly separate: (i) the algebra of values, (ii) the event structure, and
(iii) the rules governing coherence and extension.

\section{Null Evidence and Updating}

The evidence-zero boundary is where probability-theoretic conditioning
typically becomes partial or version-dependent.  Standard measure theory uses
regular conditional probabilities and disintegration to define conditional
probabilities with null evidence, with uniqueness only up to almost-sure
equivalence~\cite{chang1997}.  Category-theoretic approaches (e.g.\ Markov
categories) also treat conditioning structurally rather than as division, again
making clear which parts come from additional structure beyond the chain rule
itself~\cite{fritz2020}.

Imprecise probability develops a complementary view: at the boundary, one may
return \emph{sets} or \emph{intervals} of admissible conditionals rather than
a single number; see standard references such as Walley~\cite{walley1991} and
modern expositions~\cite{troffaes2014}.
Cred's evidence-zero underdetermination sits naturally with this perspective and
suggests an explicitly set-valued update semantics as a viable Part~3 direction.

\section{Self-Reference: Fixed Points Without Triviality}

Cred's most visible self-reference result is the forced negation fixed point
$c=\negc c \Rightarrow c=\tfrac12$.  This is algebraic and should be contrasted
with proof-theoretic fixed-point traditions:
\begin{itemize}[nosep]
\item Kripke's fixed-point theory of truth uses monotone operators and fixed
points to give consistent semantics for truth predicates~\cite{kripke1975}.
\item Revision theory (Gupta--Belnap) models truth as evolving under revisions,
yielding dynamical semantics for self-reference~\cite{gupta1993}.
\item Many-valued approaches to paradox and vagueness (including fuzzy logics)
analyze how graded truth values interact with liar-style constraints; see, e.g.,
surveys in the fuzzy-logic literature~\cite{bergmann2008}.
\end{itemize}

Part~3 should treat the Part~1 fixed point as the \emph{base case} and position
it explicitly relative to these established fixed-point traditions, rather than
as a wholesale solution to semantic paradox.

\section{Synthesis: What Cred Adds (and Refuses to Add)}

The literature map suggests a clean synthesis statement:
\begin{quote}
Cred isolates the chain rule as an algebraic constraint on conditioning and
combines it with the product De~Morgan triplet on $[0,1]$; it is compatible with
both (i) truth-functional many-valued logics (when a fixed joint is chosen) and
(ii) coherence/probability-style semantics (when joints are supplied externally).
Its distinguishing choice is to leave the evidence-zero boundary underdetermined
by default, in line with conditional-event and imprecise-probability viewpoints,
rather than forcing ex falso via residuation.
\end{quote}

This synthesis makes it clear which parts of Cred are inherited (t-norm logic,
primitive conditional probability, conditional events/coherence) and which parts
are a deliberate stance (evidence-zero underdetermination as a first-class
feature rather than an exception to be patched).

\section{Part~3 Roadmap (Paper-Grade Targets)}

The goal of Part~3 is to turn the above synthesis into concrete, checkable
results and a clear workflow for using Cred:
\begin{enumerate}[label=(\arabic*),nosep]
\item \textbf{A graded proof theory layer:} specify a minimal proof calculus
for a chosen consequence relation (e.g.\ a threshold semantics on $[0,1]$ or a
designated-values semantics after collapse) and prove soundness with respect to
the Part~2 semantics.
\item \textbf{A null-evidence update semantics:} formalize at least one update
mechanism that extends Part~2 updates to the evidence-zero boundary (likely
set-valued/interval-valued), and relate it to coherence and imprecise probability.
\item \textbf{A self-reference interface:} state precisely what the Part~1 fixed
point result does and does not address, with explicit comparison to Kripke and
revision-theoretic fixed points.
\item \textbf{A graded-foundations worked example:} pick a concrete fragment
(e.g.\ graded predicates over a finite domain, or a tiny graded set fragment) and
work through extensionality/comprehension choices, using the many-valued set
theory literature to avoid reinventing known pitfalls.
\end{enumerate}

\begin{thebibliography}{99}

\bibitem{adams1975}
Adams, E.W.
\newblock \emph{The Logic of Conditionals}.
\newblock Reidel, 1975.

\bibitem{bergmann2008}
Bergmann, M.
\newblock \emph{An Introduction to Many-Valued and Fuzzy Logic}.
\newblock Cambridge University Press, 2008.

\bibitem{calabrese1987}
Calabrese, P.
\newblock An algebraic synthesis of the foundations of logic and probability.
\newblock \emph{Information Sciences}, 42(2):187--237, 1987.

\bibitem{chang1997}
Chang, J.T. and Pollard, D.
\newblock Conditioning as disintegration.
\newblock \emph{Statistica Neerlandica}, 51(3):287--317, 1997.

\bibitem{cintula2011}
Cintula, P., H\'ajek, P., and Noguera, C., editors.
\newblock \emph{Handbook of Mathematical Fuzzy Logic}, Vols.~1--2.
\newblock College Publications, 2011.

\bibitem{coletti2002}
Coletti, G. and Scozzafava, R.
\newblock \emph{Probabilistic Logic in a Coherent Setting}.
\newblock Trends in Logic 15. Kluwer/Springer, 2002.

\bibitem{definetti1936}
de Finetti, B.
\newblock La logique de la probabilit\'e.
\newblock \emph{Actes du Congr\`es International de Philosophie Scientifique},
4:31--39, 1936.

\bibitem{definetti1937}
de Finetti, B.
\newblock \emph{La pr\'evision: ses lois logiques, ses sources subjectives}.
\newblock Annales de l'Institut Henri Poincar\'e, 7(1):1--68, 1937.

\bibitem{egre2021}
\'Egr\'e, P., Rossi, L., and Sprenger, J.
\newblock De Finettian logics of indicative conditionals {Part II}: Proof theory and algebraic semantics.
\newblock \emph{Journal of Philosophical Logic}, 50:215--247, 2021.

\bibitem{fritz2020}
Fritz, T.
\newblock A synthetic approach to {Markov} kernels, conditional independence and theorems on sufficient statistics.
\newblock \emph{Advances in Mathematics}, 370:107239, 2020.

\bibitem{gilio2013}
Gilio, A. and Sanfilippo, G.
\newblock Conjunction, disjunction and iterated conditioning of conditional events.
\newblock In Kruse, R. et al., editors, \emph{Synergies of Soft Computing and Statistics},
pages 399--407. Springer, 2013.

\bibitem{gilio2014}
Gilio, A. and Sanfilippo, G.
\newblock Conditional random quantities and compounds of conditionals.
\newblock \emph{Studia Logica}, 102(4):709--729, 2014.

\bibitem{goguen1969}
Goguen, J.A.
\newblock The logic of inexact concepts.
\newblock \emph{Synthese}, 19:325--373, 1969.

\bibitem{gupta1993}
Gupta, A. and Belnap, N.
\newblock \emph{The Revision Theory of Truth}.
\newblock MIT Press, 1993.

\bibitem{galatos2007}
Galatos, N., Jipsen, P., Kowalski, T., and Ono, H.
\newblock \emph{Residuated Lattices: An Algebraic Glimpse at Substructural Logics}.
\newblock Elsevier, 2007.

\bibitem{hajek1998}
H\'ajek, P.
\newblock \emph{Metamathematics of Fuzzy Logic}.
\newblock Kluwer, 1998.

\bibitem{kripke1975}
Kripke, S.
\newblock Outline of a theory of truth.
\newblock \emph{Journal of Philosophy}, 72(19):690--716, 1975.

\bibitem{klement2000}
Klement, E.P., Mesiar, R., and Pap, E.
\newblock \emph{Triangular Norms}.
\newblock Kluwer, 2000.

\bibitem{lewis1976}
Lewis, D.
\newblock Probabilities of conditionals and conditional probabilities.
\newblock \emph{Philosophical Review}, 85(3):297--315, 1976.

\bibitem{popper1959}
Popper, K.R.
\newblock \emph{The Logic of Scientific Discovery}.
\newblock Hutchinson, London, 1959.

\bibitem{pavelka1979i}
Pavelka, J.
\newblock On fuzzy logic {I}: Many-valued rules of inference.
\newblock \emph{Zeitschrift f\"ur mathematische Logik und Grundlagen der Mathematik} (now \emph{Mathematical Logic Quarterly}), 25:45--52, 1979.

\bibitem{pavelka1979ii}
Pavelka, J.
\newblock On fuzzy logic {II}: Enriched residuated lattices and semantics of propositional calculi.
\newblock \emph{Zeitschrift f\"ur mathematische Logik und Grundlagen der Mathematik} (now \emph{Mathematical Logic Quarterly}), 25:119--134, 1979.

\bibitem{pavelka1979iii}
Pavelka, J.
\newblock On fuzzy logic {III}: Semantical completeness of some many-valued propositional calculi.
\newblock \emph{Zeitschrift f\"ur mathematische Logik und Grundlagen der Mathematik} (now \emph{Mathematical Logic Quarterly}), 25:447--464, 1979.

\bibitem{renyi1955}
R\'enyi, A.
\newblock On a new axiomatic theory of probability.
\newblock \emph{Acta Mathematica Academiae Scientiarum Hungaricae}, 6:285--335, 1955.

\bibitem{takeuti1984}
Takeuti, G. and Titani, S.
\newblock Fuzzy logic and fuzzy set theory.
\newblock \emph{Journal of Symbolic Logic}, 49(3):851--866, 1984.

\bibitem{troffaes2014}
Troffaes, M.C.M. and de Cooman, G.
\newblock \emph{Lower Previsions}.
\newblock Wiley, 2014.

\bibitem{vanfraassen1976}
van Fraassen, B.C.
\newblock Probabilities of conditionals.
\newblock In Harper, W.L. and Hooker, C.A., editors, \emph{Foundations of Probability Theory, Statistical Inference, and Statistical Theories of Science},
Vol.~1, pages 261--308. Reidel, 1976.

\bibitem{vanfraassen1983}
van Fraassen, B.C.
\newblock Gentlemen's wagers: Relevant logic and probability.
\newblock \emph{Philosophical Studies}, 43:47--61, 1983.

\bibitem{walley1991}
Walley, P.
\newblock \emph{Statistical Reasoning with Imprecise Probabilities}.
\newblock Chapman and Hall, 1991.

\bibitem{zadeh1965}
Zadeh, L.A.
\newblock Fuzzy sets.
\newblock \emph{Information and Control}, 8:338--353, 1965.

\end{thebibliography}

\end{document}
